\documentclass[blue]{beamer}
\usetheme{metropolis}
\usepackage{amssymb,latexsym}
\usepackage{amsmath}
\usepackage{amsthm}
\usepackage{graphicx}
\usepackage{subfigure}
\usepackage{hyperref}

\definecolor{Red}{rgb}{1,0,0}
\definecolor{Blue}{rgb}{0,0,1}
\definecolor{darkblue}{rgb}{0,0,.75}
\definecolor{Green}{rgb}{0,1,0}
\definecolor{magenta}{rgb}{1,0,.6}
\definecolor{lightblue}{rgb}{0,.5,1}
\definecolor{lightpurple}{rgb}{.6,.4,1}
\definecolor{gold}{rgb}{.6,.5,0}
\definecolor{orange}{rgb}{1,0.4,0}
\definecolor{hotpink}{rgb}{1,0,0.5}
\definecolor{newcolor2}{rgb}{.5,.3,.5}
\definecolor{newcolor}{rgb}{0,.3,1}
\definecolor{newcolor3}{rgb}{1,0,.35}
\definecolor{darkgreen1}{rgb}{0, .35, 0}
\definecolor{darkgreen}{rgb}{0, .6, 0}
\definecolor{darkred}{rgb}{.75,0,0}
%\usepackage{beamerthemesplit}
\usepackage{color}
\usepackage{multirow}

\usepackage{lipsum}
\usefonttheme{professionalfonts}
\newcommand\blfootnote[1]{%
  \begingroup
  \renewcommand\thefootnote{}\footnote{#1}%
  \addtocounter{footnote}{-1}%
  \endgroup
}
\newcommand{\E}{\mathbb{E}}
\newcommand{\bi}{\begin{itemize}}
\newcommand{\ei}{\end{itemize}}
%\AtBeginSection{}
\setbeamertemplate{section in toc}[sections numbered]
\setbeamerfont{itemize/enumerate body}{size=\normalsize}
\setbeamerfont{itemize/enumerate subbody}{size=\normalsize}

%\usepackage{amsmath,amssymb,amsthm}
%\usepackage{graphicx}
\usepackage{booktabs}
\usepackage{tikz}
\usetikzlibrary{arrows.meta,positioning,decorations.pathreplacing}
%\usepackage{hyperref}



\title[Chapter22]{Chapter 22: Heterogeneous Firms}
\author[Toshihiko Mukoyama]{Chapter author: Toshihiko Mukoyama}

\date[April 2026]{April 2026}





\begin{document}
  \frame
  {
    \titlepage
  }


% ============================================================
% OUTLINE
% ============================================================
\begin{frame}{Outline}
	\tableofcontents
\end{frame}

% ============================================================
\section{Introduction}
% ============================================================

\begin{frame}{Motivation}
	Standard macroeconomic models assume an \textbf{aggregate production function}:
	\[
	Y = F(K, L)
	\]
	\begin{itemize}
		\item Justified if all firms have homogeneous production functions
		\item In reality, firms are heterogeneous in many dimensions:
		\begin{itemize}
			\item Large vs.\ small
			\item Young vs.\ old
			\item Growing vs.\ contracting
			\item Productive vs.\ unproductive
			\item Differentiated products
		\end{itemize}
	\end{itemize}
\end{frame}

\begin{frame}{Why Does Firm Heterogeneity Matter?}
	Key policy and research questions that require explicit firm heterogeneity:
	\begin{itemize}
		\item How should we encourage (or discourage) the \textbf{entry of new firms}?
		\item Should we support \textbf{growing firms}, and if so how?
		\item Should we be concerned about the growing prominence of \textbf{``mega-firms''}?
		\item What are the causes and consequences of the recent \textbf{decline in business dynamism}?
		\item Does the representative-firm assumption give \textbf{misleading answers} to standard policy questions?
	\end{itemize}
\end{frame}

\begin{frame}{Chapter Roadmap}
	\begin{enumerate}
		\item A simple model showing how heterogeneity affects aggregation
		\item Empirical facts on U.S.\ firm heterogeneity
		\item Reallocation and misallocation
		\item General equilibrium: Hopenhayn-Rogerson (1993) framework
		\item Alternative market arrangements (monopolistic competition, oligopoly)
		\item Business cycles and heterogeneous firms
		\item Endogenous productivity: Klette-Kortum (2004) model
	\end{enumerate}
\end{frame}

% ============================================================
\section{A Simple Model}
% ============================================================

\begin{frame}{Setup}
	\begin{itemize}
		\item Unit mass of firms, $i \in [0,1]$
		\item Homogeneous good, perfect competition
		\item Production function of firm $i$:
		\[
		y_i = a_i F(\mathbf{x}_i)^\gamma, \qquad \gamma \in (0,1)
		\]
		\item $a_i$: firm-specific productivity (heterogeneous across firms)
		\item $\mathbf{x}_i$: input vector for firm $i$
		\item $F(\mathbf{x}_i)$: constant returns to scale
		\item Overall production: \textbf{decreasing returns to scale} (because $\gamma < 1$)
	\end{itemize}
\end{frame}

\begin{frame}{Why Decreasing Returns?}
	\textbf{Key modeling assumption}: $\gamma < 1$ ensures DRS at the firm level
	\medskip
	
	With constant returns ($\gamma = 1$):
	\begin{itemize}
		\item The most productive firm(s) take over the entire economy
		\item Either: monopoly/oligopoly $\Rightarrow$ contradicts perfect competition
		\item Or: only the most efficient firms operate $\Rightarrow$ replicates homogeneous case
	\end{itemize}
	\medskip
	With $\gamma < 1$:
	\begin{itemize}
		\item Firms of different productivities can coexist
		\item Each firm has an optimal finite scale
		\item Richer implications for aggregation
	\end{itemize}
\end{frame}

\begin{frame}{Two-Step Optimization}
	\textbf{Step 1}: Cost minimization (common across firms)
	\[
	\min_{\mathbf{x}} \; \mathbf{px} \quad \text{s.t.} \quad F(\mathbf{x}) = 1
	\]
	Let $\mathbf{x}^*$ be the solution, $c \equiv \mathbf{px}^*$ the unit cost.
	
	\bigskip
	\textbf{Step 2}: Optimal scale for firm $i$
	
	Let $m_i = F(\mathbf{x}_i)$ be combined input. Then:
	\[
	\max_{m_i} \; a_i m_i^\gamma - c m_i
	\]
	
	First-order condition:
	\[
	a_i m_i^{\gamma - 1} = \frac{c}{\gamma}
	\]
\end{frame}

\begin{frame}{Aggregation}
	From the FOC: $y_i = (c/\gamma) m_i$ for all $i$. Summing:
	\[
	Y = \frac{c}{\gamma} M
	\]
	where $Y = \int y_i \, di$ is total output and $M = \int m_i \, di = F(\mathbf{X})$.
	
	\bigskip
	Define \textbf{aggregate productivity}:
	\[
	A \equiv \left( \int a_i^{\frac{1}{1-\gamma}} \, di \right)^{1-\gamma}
	\]
	
	Then the \textbf{aggregate production function} is:
	\[
	Y = A F(\mathbf{X})^\gamma
	\]
	Heterogeneity matters through $A$---the distribution of $a_i$ shapes aggregate output.
\end{frame}

\begin{frame}{Lognormal Example}
	Suppose $\log(a_i) \sim N(\nu - \sigma^2/2, \; \sigma^2)$
	
	\begin{itemize}
		\item Average productivity: $\int a_i \, di = \exp(\nu)$
		\item Aggregate productivity:
	\end{itemize}
	\[
	{A = \exp\left( \nu + \frac{\gamma}{1-\gamma} \frac{1}{2} \sigma^2 \right)}
	\]
	
	\bigskip
	\textbf{Key insight}: Even holding the average of $a_i$ fixed, a higher dispersion $\sigma$ raises $A$.
	
	\bigskip
	\textbf{Intuition}: Productive resources are endogenously allocated---productive firms use more input and have a greater presence in aggregate production. More dispersion $\Rightarrow$ more room to allocate to productive firms in the right tail.
\end{frame}

\begin{frame}{Two Fundamental Questions}
	Throughout this chapter, we keep two questions in mind:
	
	\bigskip
	\begin{enumerate}
		\item \textbf{How is the distribution of $a_i$ determined?}
		\begin{itemize}
			\item Exogenous vs.\ endogenous productivity
			\item Entry, exit, and innovation
		\end{itemize}
		
		\bigskip
		\item \textbf{How does the economy allocate resources to different firms?}
		\begin{itemize}
			\item Efficient allocation vs.\ misallocation
			\item Frictions, distortions, and policy
		\end{itemize}
	\end{enumerate}
\end{frame}


% ============================================================
\section{Firm Heterogeneity in the Data}
% ============================================================

\begin{frame}{Data Source}
	\begin{itemize}
		\item All figures from the U.S.\ Census Bureau's \textbf{Business Dynamics Statistics (BDS)}
		\item Covers the universe of U.S.\ employer establishments and firms
		\item Key distinction:
		\begin{itemize}
			\item \textbf{Establishment}: a fixed physical location of economic activity
			\item \textbf{Firm}: a collection of establishments under common ownership
		\end{itemize}
		\item Over 5 million firms and 7 million establishments in the U.S.
	\end{itemize}
\end{frame}

\begin{frame}{Firm Size Distribution}
	
	Distribution of firm size in 2022
	\begin{figure}[h]
		\centering
		\includegraphics[scale=0.5]{FigsChapter22/fig1_firm_size_dist}
	\end{figure}
\vspace{-20pt}
	\begin{itemize}
		\item Highly dispersed distribution
		\item Majority are very small firms (1--4 employees)
		\item Over 10,000 firms with 1,000+ employees
		\item Over 1,000 firms with 10,000+ employees
	\end{itemize}
\end{frame}



\begin{frame}{Employment Shares}
	Employment share of each size category in 2022
	\begin{figure}[h]
	\centering
	\includegraphics[scale=0.5]{FigsChapter22/fig2_emp_share_firm_size_dist}
\end{figure}
\vspace{-20pt}
	\begin{itemize}
		\item Small firms are numerous but large firms employ most workers
		\item $\approx 50\%$ of employees work at firms with $1{,}000+$ employees
		\item $\approx 30\%$ work at very large firms ($10{,}000+$ employees)
	\end{itemize}
\end{frame}


\begin{frame}{Establishment Size Distribution}
	
	Distribution of establishment size in 2022
\begin{figure}[h]
	\centering
	\includegraphics[scale=0.5]{FigsChapter22/fig3_estab_size_dist}
\end{figure}
\vspace{-20pt}
	\begin{itemize}
		\item $\approx 50\%$ are very small (1--4 employees)

	\end{itemize}
\end{frame}


\begin{frame}{Firm Size and Establishment Share}
	Fraction of establishments owned by each firm size category in 2022
	\begin{figure}[h]
		\centering
		\includegraphics[scale=0.5]{FigsChapter22/fig4_estabs_share_firm_size_dist}
	\end{figure}
	\vspace{-20pt}
	\begin{itemize}

		\item Large firms own many establishments
		\item Firms with $1{,}000+$ employees own avg $\approx 100$ establishments 
		\item Firms with $10{,}000+$ employees own avg $\approx 600$ establishments
	\end{itemize}
\end{frame}
\begin{frame}{Job Creation and Destruction}
	\textbf{Definitions}: Let $\ell_{it}$ be employment of establishment $i$ at year $t$
	\[
	JC_t \equiv \frac{\sum_{i:\ell_{it}>\ell_{i,t-1}} (\ell_{it} - \ell_{i,t-1})}{\bar{L}_t}, \qquad
	JD_t \equiv \frac{\sum_{i:\ell_{it}<\ell_{i,t-1}} (\ell_{i,t-1} - \ell_{it})}{\bar{L}_t}
	\]
	where $\bar{L}_t = (L_t + L_{t-1})/2$.
	
	\bigskip
	Annual job creation and destruction rates (establishment level)
	\vspace{-10pt}
	\begin{figure}[h]
		\centering
		\includegraphics[scale=0.32]{FigsChapter22/fig5_JCrate_JDrate}
	\end{figure}
\end{frame}


\begin{frame}{Properties of Job Flows}
	Three notable properties:
	
	\bigskip
	\begin{enumerate}
		\item \textbf{Large magnitudes}: Both JC and JD rates exceed 10\% in any given year
		
		\bigskip
		\item \textbf{Cyclical}: In recessions, JC declines and JD increases
		
		\bigskip
		\item \textbf{Declining trend}: Both rates have declined over time $\Rightarrow$ \textbf{``declining business dynamism''}
	\end{enumerate}
\end{frame}

\begin{frame}{Entry and Exit Rates}
	Annual establishment entry and exit rates
	\begin{figure}[h]
	\centering
	\includegraphics[scale=0.3]{FigsChapter22/fig6_estab_entry_exit_rates}
\end{figure}
\vspace{-10pt}	
	Similar properties as JC/JD:
	\begin{itemize}
		\item Large rates (entry $\approx 10$--$14\%$, exit $\approx 8$--$11\%$)
		\item Cyclical patterns
		\item Declining trend over recent decades
	\end{itemize}
\end{frame}

\begin{frame}{Rise of Large Firms}
	Fraction of employees working in firms with 10,000+ employees
	\begin{figure}[h]
	\centering
	\includegraphics[scale=0.3]{FigsChapter22/fig7_emp_share_10000+}
\end{figure}
\vspace{-10pt}
	\begin{itemize}
		\item Steady increase since early 1990s
		\item Large firms are increasingly dominating the U.S.\ economy
		\item Raises concerns about \textbf{market power}
	\end{itemize}
\end{frame}



\begin{frame}{Rising Markups}
	Markup of U.S.\ public firms
\begin{figure}[h]
	\centering
	\includegraphics[scale=0.3]{FigsChapter22/fig8_markup}
\end{figure}
\vspace{-15pt}
	\begin{itemize}
		\item Average markup (price/marginal cost) of U.S.\ public firms
		\item Increasing trend since the 1980s
		\item Consistent with rising dominance of large firms
		\item Active research area; see Miller (2025) and Syverson (2025) for surveys
	\end{itemize}
\end{frame}

% ============================================================
\section{Reallocation and Misallocation}
% ============================================================

\begin{frame}{How Much Does Reallocation Matter?}
	\textbf{Foster, Haltiwanger, and Krizan (2001)} decomposition:
	
	\medskip
	Define output-weighted average productivity: $\bar{A}_t \equiv \sum_i s_{it} a_{it}$
	
	where $s_{it}$ is the output share of establishment $i$.
	
	\bigskip
	\[
	\Delta \bar{A}_t = \underbrace{\sum_{i \in C} s_{i,t-1} \Delta a_{it}}_{\text{within}} + \underbrace{\sum_{i \in C} (a_{i,t-1} - \bar{A}_{t-1}) \Delta s_{it}}_{\text{between}} + \underbrace{\sum_{i \in C} \Delta a_{it} \Delta s_{it}}_{\text{cross}}
	\]
	\[
	+ \underbrace{\sum_{i \in N} s_{it}(a_{it} - \bar{A}_{t-1})}_{\text{entry}} - \underbrace{\sum_{i \in X} s_{i,t-1}(a_{i,t-1} - \bar{A}_{t-1})}_{\text{exit}}
	\]
	
	\medskip
	$C$: continuing, $N$: new entrants, $X$: exiting establishments.
\end{frame}

\begin{frame}{Decomposition Results}
	Five sources of aggregate productivity growth:
	\begin{enumerate}
		\item \textbf{Within}: each establishment raises its own productivity
		\item \textbf{Between}: high-productivity establishments gain market share
		\item \textbf{Cross}: interaction of productivity growth and share changes
		\item \textbf{Entry}: entrants are above average
		\item \textbf{Exit}: exiting establishments are below average
	\end{enumerate}
	
	\bigskip
	\textbf{Key finding} (U.S.\ manufacturing, 1977--1987):
	\begin{itemize}
		\item Only ${45\%}$ from the within effect
		\item The remaining ${55\%}$ from reallocation (between, cross, entry, exit)
	\end{itemize}
	
	$\Rightarrow$ \textbf{Reallocation is crucial for aggregate productivity growth}
\end{frame}
\begin{frame}{Misallocation: Setup}
	\textbf{Restuccia and Rogerson (2008), Hsieh and Klenow (2009)}: firm-specific distortions cause misallocation.
	
	\bigskip
	Add firm-specific output tax $\tau_i$ to the simple model:
	\[
	\max_{m_i} \; (1-\tau_i) a_i m_i^\gamma - c m_i
	\]
	
	\medskip
	GDP is still $Y = \int y_i  di$. Aggregate production function: $Y = A F(\mathbf{X})^\gamma$
	
	\medskip
	But now:
	\[
	A = \frac{\left[ \int a_i^{\frac{1}{1-\gamma}} (1-\tau_i)^{\frac{\gamma}{1-\gamma}} \, di \right]^{1-\gamma}}{\left[ \int a_i^{\frac{1}{1-\gamma}} (1-\tau_i)^{\frac{1}{1-\gamma}} \, di \right]^{\gamma}}
	\]
	
	Reduces to the undistorted $A$ when $\tau_i = 0$ for all $i$.
\end{frame}

\begin{frame}{Misallocation: Lognormal Case}
	Assume $(a_i, 1-\tau_i)$ are bivariate lognormal:
	\[
	(\log a_i, \log(1-\tau_i)) \sim N(\mu, \Sigma)
	\]
	with $\mu = (\nu_a - \sigma_a^2/2, \; \nu_\tau - \sigma_\tau^2/2)$ and $\text{Var}(\log(1-\tau_i)) = \sigma_\tau^2$.
	
	\bigskip
	\textbf{Result}:
	\[
	{A = \exp\left( \nu_a + \frac{\gamma}{1-\gamma} \frac{1}{2}(\sigma_a^2 - \sigma_\tau^2) \right)}
	\]
	
	\bigskip
	\textbf{Key insights}:
	\begin{itemize}
		\item Higher dispersion in taxes $\sigma_\tau$ \textbf{reduces} aggregate productivity
		\item Productive firms may not expand if $(1-\tau_i) a_i$ is low
		\item Unproductive firms may over-expand if $(1-\tau_i) a_i$ is high
		\item The {level} of the tax $\nu_\tau$ does not affect $A$ (no distortion in allocation)
	\end{itemize}
\end{frame}

\begin{frame}{Sources of Misallocation}
	\textbf{Specific policies and institutions} that cause misallocation:
	\begin{itemize}
		\item Size-specific taxes and regulations
		\item Entry regulations and barriers
		\item Regulations governing hiring and firing
		\item Financial frictions (credit constraints)
		\item Trade barriers
	\end{itemize}
	
	\bigskip
	\textbf{Quantitative importance}: Hsieh and Klenow (2009) estimate that reallocating resources to equalize marginal products would raise manufacturing TFP by:
	\begin{itemize}
		\item 30--50\% in China
		\item 40--60\% in India
	\end{itemize}
\end{frame}


% ============================================================
\section{Firm Heterogeneity in General Equilibrium}
% ============================================================

\begin{frame}{The Hopenhayn-Rogerson (1993) Framework}
	\textbf{Goal}: Analyze reallocation frictions in general equilibrium with forward-looking firms.
	
	\bigskip
	\textbf{Key features}:
	\begin{itemize}
		\item Continuum of firms with idiosyncratic productivity shocks
		\item Endogenous entry and exit
		\item Labor as the only input
		\item Free entry condition
		\item Representative consumer
		\item Focus on steady state
	\end{itemize}
	
	\bigskip
	\textbf{Experiments}: firing taxes and entry barriers (Moscoso Boedo and Mukoyama, 2012).
\end{frame}

\begin{frame}{Production and Profits}
	\begin{itemize}
		\item Production function: $y_t = a_t \ell_t^\gamma$
		\item $a_t$: idiosyncratic productivity (stochastic)
		\item Firms pay:
		\begin{itemize}
			\item Wages: $w_t \ell_t$
			\item Fixed operation cost: $c_f$ (units of goods, per period)
			\item Firing taxes: $\tau \max(0, \ell_{t-1} - \ell_t)$
		\end{itemize}
	\end{itemize}
	
	\bigskip
	\textbf{Flow profit}:
	\[
	\pi(\ell_{t-1}, \ell_t, a_t) = a_t \ell_t^\gamma - w_t \ell_t - c_f - \tau \max(0, \ell_{t-1} - \ell_t)
	\]
	
	\begin{itemize}
		\item Output price normalized to 1
		\item Firing taxes transferred back to consumers
	\end{itemize}
\end{frame}

\begin{frame}{Timing and Productivity Process}
	\textbf{Within-period timing}:
	\begin{enumerate}
		\item Incumbent decides whether to \textbf{exit}
		\begin{itemize}
			\item If exit: pay firing cost $\tau \ell_{t-1}$
		\end{itemize}
		\item If stay: receive current $a_t$ from
		\[
		\log(a_t) = \alpha + \rho \log(a_{t-1}) + \varepsilon_t, \qquad \varepsilon_t \sim N(0, \sigma^2)
		\]
		\item Choose employment $\ell_t$, produce
	\end{enumerate}
	
	\bigskip
	\textbf{Dynamic element}: With $\tau > 0$, hiring is dynamic---firm foresees future firing costs $\Rightarrow$ reluctance to hire even with positive shocks.
\end{frame}

\begin{frame}{Firm's Dynamic Program}
	\[
	W(a, \ell_{-1}) = \max_{\ell} \; \pi(\ell_{-1}, \ell, a) + \beta \max\left\{ \E[W(a', \ell) \mid a], \; -\tau \ell \right\}
	\]
	
	\begin{itemize}
		\item $\beta \in (0,1)$: consumer's discount factor
		\item Inner $\max$: \textbf{exit decision}
		\begin{itemize}
			\item Continue: get $\E[W(a', \ell) \mid a]$
			\item Exit: pay firing cost $-\tau \ell$
		\end{itemize}
	\end{itemize}
\end{frame}

\begin{frame}{Entry and Consumer}
	\textbf{Free entry}:
	\[
	W^e = c_e + \kappa
	\]
	where $c_e$ = technological entry cost, $\kappa$ = policy-related entry barriers.
	\[
	W^e = \int (W(a, 0) + c_f) \, d\nu(a)
	\]
	$\nu(a)$: exogenous distribution for entrant productivity.
	
	\bigskip
	\textbf{Representative consumer}:
	\[
	\sum_{t=0}^\infty \beta^t [u(C_t) - \chi L_t^s]
	\]
	\textbf{Steady-state FOC}: $w u'(wL^s + \Pi + R) = \chi$
	
	\medskip
	where $\Pi$ = total firm profits, $R$ = total transfers.
\end{frame}

\begin{frame}{Block-Recursive Structure}
	\textbf{Key computational insight} (Kaas, 2021):
	
	\bigskip
	\begin{enumerate}
		\item Solve $W(a, \ell_{-1})$ as function of $w$ ($\Rightarrow$ get $W^e(w)$)
		\item Free entry pins down $w^*$ ($\Rightarrow$ $W^e(w^*) = c_e + \kappa$)
		\item Given $w^*$ and entry mass $M$: compute stationary distribution
		\item Compute $L^d(M)$, $\Pi(M)$, $R(M)$
		\item Labor market clearing $L^d(M) = L^s(\Pi(M), R(M))$ pins down $M$
	\end{enumerate}
	
	\bigskip
	\textbf{Advantage}: No need to track firm distribution to solve for prices (``block recursive structure")
	
	\medskip
	$\Rightarrow$ Much easier computationally than heterogeneous-consumer models. (Lee and Mukoyama, 2018)
\end{frame}

\begin{frame}{Calibration}
	\begin{itemize}
		\item One period = 5 years
		\item $u(c) = \log(c)$
		\item $\beta = 0.8$ (4\% annual discount rate)
		\item $\gamma = 0.64$ (labor share)
		\item Baseline: $(\tau, \kappa) = (0, 0)$ calibrated to the US economy
	\end{itemize}
	
	\bigskip
	\textbf{Targets}:
	\begin{itemize}
		\item $\alpha = 0.076$: average firm size $\approx 62$ employees
		\item $\rho = 0.93$: autocorrelation of $\log(\ell)$ is 0.93
		\item $\sigma = 0.253$: variance of employment growth is 0.53
		\item $c_f = 18.0$: exit rate $\approx 37\%$ (over 5 years)
		\item $c_e = 9.04$: free entry holds with $w = 1$
	\end{itemize}
\end{frame}

\begin{frame}{Effects of Firing Taxes}
	\begin{table}
		\centering
		\begin{tabular}{lccc}
			\toprule
			& $\tau = 0$ & $\tau = 0.1$ & $\tau = 0.2$ \\
			\midrule
			Wage & 1.000 & 0.977 & 0.957 \\
			Total output & 100 & 97.7 & 95.7 \\
			Total employment & 100 & 98.3 & 97.4 \\
			Labor productivity & 100 & 99.4 & 98.3 \\
			JC (= JD) rate & 0.28 & 0.25 & 0.21 \\
			\bottomrule
		\end{tabular}
		\caption{Model results with firing taxes ($\kappa = 0$).}
	\end{table}
	
	\begin{itemize}
		\item $\tau = 0.1$ $\approx$ 6 months' salary
		\item $\tau = 0.2$ $\approx$ 1 year's salary
	\end{itemize}
\end{frame}

\begin{frame}{Interpreting the Firing Tax Results}
	\begin{enumerate}
		\item \textbf{Employment declines}: Hiring effect dominates firing effect
		\begin{itemize}
			\item Firms fire less $\Rightarrow$ pushes $L$ up
			\item Firms hire less (foresee future firing costs) $\Rightarrow$ pushes $L$ down
			\item Net: employment falls
		\end{itemize}
		
		\bigskip
		\item \textbf{Labor productivity ($Y/L$) declines}: Misallocation
		\begin{itemize}
			\item Good-$a$ firms don't expand enough
			\item Bad-$a$ firms don't contract enough
			\item Marginal products of labor become dispersed
		\end{itemize}
		
		\bigskip
		\item \textbf{JC rate declines}: Less reallocation
		\begin{itemize}
			\item Reluctance to hire and fire
			\item Closely linked to productivity loss from misallocation
		\end{itemize}
	\end{enumerate}
\end{frame}

\begin{frame}{Effects of Entry Barriers}
	\begin{table}
		\centering
		\begin{tabular}{lccc}
			\toprule
			& $\kappa = 0$ & $\kappa = 0.5$ & $\kappa = 5.0$ \\
			\midrule
			Wage & 1.000 & 0.986 & 0.879 \\
			Total output & 100 & 98.6 & 87.9 \\
			Total employment & 100 & 99.5 & 96.2 \\
			Labor productivity & 100 & 99.1 & 91.4 \\
			JC (= JD) rate & 0.28 & 0.28 & 0.28 \\
			\bottomrule
		\end{tabular}
		\caption{Model results with entry barriers ($\tau = 0$).}
	\end{table}
	
	\begin{itemize}
		\item $\kappa = 5.0$: comparable to low-income country entry costs 
		\item Substantial output and productivity losses
	\end{itemize}
\end{frame}

\begin{frame}{Interpreting the Entry Barrier Results}
	\textbf{Mechanism}: Higher $\kappa$ $\Rightarrow$ higher firm value in equilibrium $\Rightarrow$ lower wage
	
	\bigskip
	Three channels through which low wage affects productivity:
	\begin{enumerate}
		\item Low-productivity firms \textbf{less likely to exit} $\Rightarrow$ drags down $A$ \quad {\color{red}$\downarrow$}
		\item Low exit rate $\Rightarrow$ low entry rate $\Rightarrow$ fewer (less productive) entrants $\Rightarrow$ raises $A$ \quad {\color{blue}$\uparrow$}
		\item Larger incumbent size $\Rightarrow$ DRS $\Rightarrow$ lower productivity \quad {\color{red}$\downarrow$}
	\end{enumerate}
	
	\bigskip
	Net effect: productivity \textbf{declines} (channels 1 \& 3 dominate).
	
	\bigskip
	\textbf{Political economy}: Entry barriers \textbf{harm entrants} but \textbf{benefit incumbents} $\Rightarrow$ conflict of interest (Mukoyama and Popov, 2014).
\end{frame}

% ============================================================
\section{Alternative Market Arrangements}
% ============================================================

\begin{frame}{Moving Beyond Perfect Competition}
	\textbf{Two takeaways}:
	\begin{enumerate}
		\item Misallocation insights carry over with minor modifications
		\item Market power enables analysis of firm heterogeneity and \textbf{markups}
	\end{enumerate}
	
	\bigskip
	We consider two alternatives:
	\begin{itemize}
		\item \textbf{Monopolistic competition} (constant markups)
		\item \textbf{Oligopoly} (endogenous markups) --- Atkeson and Burstein (2008)
	\end{itemize}
\end{frame}

\begin{frame}{Monopolistic Competition: Setup}
	\textbf{Final good} (CES aggregation / Dixit--Stiglitz):
	\[
	Y = \left( \int y_i^{\frac{\sigma-1}{\sigma}} \, di \right)^{\frac{\sigma}{\sigma-1}}, \qquad \sigma > 1
	\]
	
	\begin{itemize}
		\item Each intermediate good $i$ produced by a monopolist
		\item Same input structure as the simple model
		\item Input markets: perfectly competitive
	\end{itemize}
	
	\bigskip
	\textbf{Inverse demand} for good $i$:
	\[
	p_i = y_i^{-1/\sigma} Y^{1/\sigma}
	\]
\end{frame}

\begin{frame}{Monopolistic Competition: Firm Problem}
	Producer $i$ solves:
	\[
	\max_{m_i} \; (a_i m_i^\gamma)^{-1/\sigma} Y^{1/\sigma} \cdot a_i m_i^\gamma - c m_i
	\]
	taking $Y$ as given (firm is small relative to aggregate).
	
	\bigskip
	\textbf{Key result} --- Constant markup pricing:
	\[
	{p_i = \frac{\sigma}{\sigma - 1} \cdot \mathcal{M}}
	\]
	where $\mathcal{M}$ is marginal cost. The markup $\sigma/(\sigma-1)$ is \textbf{constant}.
	
	\bigskip
	\textbf{Aggregate production function}: $Y = A F(\mathbf{X})^\gamma$ with modified $A$:
	\[
	A \equiv \left( \int a_i^{\frac{\sigma}{\sigma-1} - \gamma} \, di \right)^{\frac{\sigma-1}{\sigma-1-\gamma}}
	\]
\end{frame}

\begin{frame}{Monopolistic Competition: Key Properties}
	\begin{itemize}
		\item Aggregate production function same form as before
		\item Aggregation uses $\sigma/(\sigma-1)$ instead of $1/(1-\gamma)$
		\item \textbf{Advantage}: Can accommodate \textbf{constant} or even some \textbf{increasing} returns to scale
		\begin{itemize}
			\item Monopoly power provides another force limiting firm size
		\end{itemize}
		\item \textbf{Limitation}: Constant markup $\Rightarrow$ cannot analyze endogenous markup changes
	\end{itemize}
\end{frame}

\begin{frame}{Oligopoly: Atkeson-Burstein (2008)}
	\textbf{Motivation}: Need endogenous markups to study:
	\begin{itemize}
		\item Rising markups since the 1980s
		\item How changes in competition affect markups
	\end{itemize}
	
	\bigskip
	\textbf{Structure}: Two layers of aggregation
	\begin{enumerate}
		\item \textbf{Across sectors}: CES with elasticity $\sigma$ (same as before)
		\[
		Y = \left( \int y_i^{\frac{\sigma-1}{\sigma}} \, di \right)^{\frac{\sigma}{\sigma-1}}
		\]
		\item \textbf{Within sector $i$}: $J$ firms (``brands'') with elasticity $\eta$
		\[
		y_i = \left( \sum_{j=1}^J q_{ij}^{\frac{\eta-1}{\eta}} \right)^{\frac{\eta}{\eta-1}}
		\]
	\end{enumerate}
	
	\textbf{Assumption}: $\eta > \sigma > 1$ (brands more substitutable than sectors).
\end{frame}

\begin{frame}{Endogenous Markup Formula}
	\textbf{Pricing rule} (Cournot competition):
	\[
	\hat{p}_{ij} = \frac{\varepsilon(s_{ij})}{\varepsilon(s_{ij}) - 1} \cdot \mathcal{M}
	\]
	where $s_{ij}$ is the \textbf{sales share} of firm $j$ in sector $i$:
	\[
	s_{ij} = \frac{\hat{p}_{ij} q_{ij}}{\sum_h \hat{p}_{ih} q_{ih}}
	\]
	and the effective elasticity is:
	\[
	\varepsilon(s_{ij}) = \left( \frac{1-s_{ij}}{\eta} + \frac{s_{ij}}{\sigma} \right)^{-1}
	\]
	
	\begin{itemize}
		\item $\varepsilon = \eta$ when $s_{ij} = 0$ (small firm $\Rightarrow$ low markup)
		\item $\varepsilon = \sigma$ when $s_{ij} = 1$ (monopoly $\Rightarrow$ high markup)
		\item Monotonically decreasing in $s_{ij}$
	\end{itemize}
\end{frame}

\begin{frame}{Endogenous Markups: Properties}
	\textbf{Symmetric firms}: $s_{ij} = 1/J$
	\[
	\varepsilon = \left( \frac{1}{\eta} \cdot \frac{J-1}{J} + \frac{1}{\sigma} \cdot \frac{1}{J} \right)^{-1}
	\]
	
	\begin{itemize}
		\item $J = 1$ (monopoly): $\varepsilon = \sigma$ $\Rightarrow$ same as monopolistic competition
		\item $J \to \infty$: $\varepsilon \to \eta$ $\Rightarrow$ most competitive case
		\item \textbf{Fewer firms $\Rightarrow$ higher markups}
	\end{itemize}
	
	\bigskip
	\textbf{Asymmetric firms}: firm heterogeneity in $a_{ij}$ feeds into markup heterogeneity through sales shares.
	
	\bigskip
	\textbf{Bertrand competition}: Same formula but with
	\[
	\varepsilon(s_{ij}) = \eta(1-s_{ij}) + \sigma s_{ij}
	\]
\end{frame}

% ============================================================
\section{Business Cycles and Heterogeneous Firms}
% ============================================================

\begin{frame}{Do Idiosyncratic Shocks Wash Out?}
	\textbf{Law of Large Numbers argument}:
	
	Let GDP $Y_t = \sum_{i=1}^N y_{it}$, with i.i.d.\ growth rates: $(y_{i,t+1} - y_{it})/y_{it} = \sigma \varepsilon_{i,t+1}$
	
	\bigskip
	Standard deviation of GDP growth:
	\[
	\sigma_Y = \sigma \sqrt{\sum_{i=1}^N \left(\frac{y_{it}}{Y_t}\right)^2}
	\]
	
	\bigskip
	\textbf{Equal-size firms}: $y_{it}/Y_t = 1/N$
	\[
	\sigma_Y = \frac{\sigma}{\sqrt{N}}
	\]
	
	With $N = 5{,}000{,}000$ and $\sigma \approx 15\%$: $\sigma_Y \approx 0.00007$ $\Rightarrow$ \textbf{negligible}!
\end{frame}

\begin{frame}{Aggregate Shocks and Firm Dynamics}
	\textbf{Standard approach}: Accept LLN, add aggregate shock $z_t$:
	\[
	y_{it} = z_t \, s_{it} \, \ell_{it}^\gamma
	\]
	where $z_t$ = aggregate productivity (as in RBC models).
	
	\bigskip
	\textbf{Results}:
	\begin{itemize}
		\item Modified Hopenhayn-Rogerson model replicates well:
		\begin{itemize}
			\item Cyclical JC and JD rates
			\item Cyclical entry and exit rates
		\end{itemize}
		\item Block-recursive structure $\Rightarrow$ computationally tractable
		\item Some firm-level statistics harder to match (Lee and Mukoyama, 2018)
	\end{itemize}
\end{frame}

\begin{frame}{Idiosyncratic Productivity and Aggregate: Hulten's Theorem}
	\textbf{Setup}: $N$ sectors, sector $i$ produces $y_i = a_i F(k_i, \ell_i, x_{i1}, \ldots, x_{iN})$
	
	where $x_{ij}$ = intermediate input from sector $j$.
	
	\bigskip
	\textbf{Theorem (Hulten, 1978)}: The first-order output effect of productivity changes:
	\[
\frac{dY}{Y} = \sum_i D_i \frac{da_i}{a_i}
	\]
	where $D_i$ is the \textbf{Domar weight}:
	\[
	D_i = \frac{p_i y_i}{\sum_i p_i c_i} = \frac{\text{Sales of sector } i}{\text{GDP (value added)}}
	\]
	
	\bigskip
	\textbf{Note}: Domar weights use \emph{sales} (not value added) $\Rightarrow$ can exceed 1!
	
	$\Rightarrow$ Downstream firms have amplified impact through input-output linkages.
\end{frame}

\begin{frame}{Hulten's Theorem: Intuition}
	\textbf{Two key intuitions}:
	
	\bigskip
	\begin{enumerate}
		\item \textbf{Why only $a_i$ matters} (not inputs)?
		\begin{itemize}
			\item Envelope theorem: in efficient economy, input adjustments have no first-order welfare effect
		\end{itemize}
		
		\bigskip
		\item \textbf{Why sales (not value added)?}
		\begin{itemize}
			\item With input-output networks, TFP improvement in a downstream firm also raises the value of intermediate inputs
		\end{itemize}
	\end{enumerate}
	
	\bigskip
	\textbf{Limitations} (relaxed by Baqaee and Farhi, 2019, 2020):
	\begin{itemize}
		\item Requires efficient economy (no misallocation)
		\item First-order approximation only
	\end{itemize}
\end{frame}

\begin{frame}{Hulten's Theorem: Simple Example}
	Two sectors:
	\begin{itemize}
		\item Sector 1 (consumption): $y_1 = a_1 x^{1-\gamma} \ell^\gamma$
		\item Sector 2 (intermediate): $y_2 = a_2 k$
	\end{itemize}
	Fixed $K$, $L$. Competitive equilibrium: $x = y_2$
	
	\[
	Y = a_1 (a_2 K)^{1-\gamma} L^\gamma
	\]
	\[
	\frac{dY}{Y} = \underbrace{1}_{D_1} \cdot \frac{da_1}{a_1} + \underbrace{(1-\gamma)}_{D_2} \cdot \frac{da_2}{a_2}
	\]
	
	\begin{itemize}
		\item $D_1 = 1$ (sales = GDP for downstream sector)
		\item $D_2 = 1 - \gamma$ (value added share of upstream sector)
		\item Domar weights sum to $1 + (1-\gamma) > 1$
	\end{itemize}
\end{frame}

\begin{frame}{Large Firms and ``Granular'' Fluctuations}
	\textbf{Gabaix (2011)}: Fat-tailed firm size distribution changes the LLN result.
	
	\medskip
	If firm sizes follow a \textbf{Pareto distribution}: $\Pr[y_i > x] = \chi x^{-\zeta}$ (the US data: $\zeta \approx 1$, called the Zipf's law)

	\medskip
	\textbf{Result}: Instead of $\sigma_Y \propto \sigma / \sqrt{N}$:
	\[
	\sigma_Y \sim \frac{v_\zeta \, \sigma}{\log(N)}
	\]
	\vspace{-20pt}
	\begin{itemize}
		\item $1/\log(1{,}000{,}000) \approx 0.072$ vs.\ $1/\sqrt{1{,}000{,}000} = 0.001$
		\item \textbf{Two orders of magnitude larger!}
		\item Idiosyncratic shocks to large firms can generate aggregate-scale fluctuations
	\end{itemize}
	
	\medskip
	Carvalho \& Grassi (2019): quantitative business cycle model driven by idiosyncratic shocks to large firms.
\end{frame}

\begin{frame}{Production Networks}
	\textbf{Additional amplification}: Domar weights $D_i = p_i y_i / \text{GDP}$
	
	\begin{itemize}
		\item Large downstream firms (Walmart, Amazon, GM): sales $\gg$ value added
		\item Their Domar weights can be much larger than their GDP share
		\item $\Rightarrow$ Idiosyncratic shocks amplified through input-output linkages
	\end{itemize}
	
	\bigskip
	\textbf{Beyond Hulten's theorem}:
	\begin{itemize}
		\item \textbf{Baqaee \& Farhi (2019)}: Second-order effects can be quantitatively important
		\item \textbf{Baqaee \& Farhi (2020)}: With distortions, network structure matters for first-order effects; decompose Solow residual into misallocation changes + ``pure'' technology growth
	\end{itemize}
\end{frame}

% ============================================================
\section{Endogenous Productivity: Klette-Kortum (2004)}
% ============================================================

\begin{frame}{Motivation}
	\begin{itemize}
		\item Previous sections: productivity $a_i$ is \textbf{exogenous}
		\item Natural question: What determines $a_i$?
		\item \textbf{Klette \& Kortum (2004)}: Endogenous productivity from \textbf{innovation}
		\item Combines:
		\begin{itemize}
			\item Quality-ladder growth models (Ch.\ 13)
			\item Firm dynamics (entry, exit, expansion, contraction)
		\end{itemize}
		\item Major strength: clear definition of a ``firm'' and its dynamics
	\end{itemize}
\end{frame}

\begin{frame}{Preferences}
	\textbf{Representative consumer}:
	\[
	\sum_{t=0}^\infty \beta^t \log(C_t)
	\]
	
	\textbf{Consumption aggregator} (unit elasticity of substitution):
	\[
	C_t = \exp\left( \int_0^1 \log\left( \sum_{k=-1}^{J_t(j)} q_t(j,k) \, c_t(j,k) \right) dj \right)
	\]
	
	\begin{itemize}
		\item $j \in [0,1]$: product index
		\item $k$: generation (quality rung), $J_t(j)$ = cutting-edge generation
		\item $q_t(j,k)$: quality of product $(j,k)$
		\item $c_t(j,k)$: consumption of product $(j,k)$
		\item Different generations are \textbf{perfect substitutes} (quality-adjusted)
	\end{itemize}
\end{frame}

\begin{frame}{Intratemporal Problem}
	\textbf{Equal expenditure shares} (unit elasticity):
	\[
	c_t(j,k) = \frac{E_t}{p_t(j,k)}
	\]
	for each $j$, consumer buys only the generation with lowest quality-adjusted price.
	
	\bigskip
	\textbf{Price index}: $P_t = \exp\left( \int_0^1 [\log p_t(j,k) - \log q_t(j,k)] \, dj \right)$
	
	Normalize $P_t = 1$.
	
	\bigskip
	\textbf{Intertemporal Euler equation} (balanced growth path):
	\[
	\frac{1}{C_t}=\beta(1+r)\frac{1}{C_{t+1}}\quad\Rightarrow\quad \frac{1}{1+r} = \frac{\beta}{1+\gamma}
	\]
	where $\gamma$ is the consumption growth rate.
\end{frame}

\begin{frame}{Production and Pricing}
	\textbf{Production}: 1 unit of labor $\Rightarrow$ 1 unit of any product.
	\begin{itemize}
		\item Unit cost = $w_t$ for all products
	\end{itemize}
	
	\bigskip
	\textbf{Limit pricing}: The cutting-edge producer sets:
	\[
	p_t(j, J_t(j)) = \lambda w_t
	\]
	where $\lambda > 1$ is the \textbf{quality step} (also the markup).
	
	\bigskip
	\textbf{Period profit per product line}:
	\[
	\pi_t = \left(1 - \frac{1}{\lambda}\right) C_t
	\]
	\begin{itemize}
		\item All product lines earn the same profit.
		\item Firm size = number of product lines owned
	\end{itemize}
\end{frame}

\begin{frame}{Innovation}
	\textbf{Innovation intensity}: $\eta$ per product line
	\begin{itemize}
		\item Cost: $w_t R(\eta)$ units of labor, $R(\cdot)$ increasing and convex
		\item Success probability: $\eta$
		\item If successful: gain a random product line (quality $\lambda$ times the current best of that line)
		\item Fixed total number of products $\rightarrow$ New product always taken from another firm (creative destruction)
	\end{itemize}
	
	\medskip
	Let $\mu$ = probability of \textbf{losing} a product line to another innovator.
	
	\medskip
	\textbf{Bellman equation} (per product line):
	\[
	V_t = \max_\eta \; \pi_t - w_t R(\eta) + \frac{1}{1+r}(1 + \eta - \mu) V_{t+1}
	\]
	Expected continuation: $(1+\eta - \mu)V_{t+1}$ (gain with prob $\eta$, lose with prob $\mu$).
\end{frame}

\begin{frame}{Balanced Growth Path}
	Divide both sides by $(1+\gamma)^t$, use $\beta/(1+\gamma) = 1/(1+r)$ and call $v\equiv V_t/(1+\gamma)^t$:
	\[
	v = \max_\eta \; \left(1 - \frac{1}{\lambda}\right) C_0 - R(\eta) + \beta(1 + \eta - \mu) v,
	\]
	
	\textbf{FOC}: $R'(\eta) = \beta v$
	
	\medskip
	\textbf{Free entry}: An entrant pays $c_e$ units of labor, gets one product line:
	\[
	v = c_e
	\]
	
	\textbf{Creative destruction}: $\mu = \eta + \nu$ \quad (incumbent innovation + entry rate)
	
	\medskip
	\textbf{Labor market clearing}:
	\[
	\frac{C_0}{\lambda} + R(\eta) + \nu c_e = L
	\]
\end{frame}

\begin{frame}{Equilibrium}
	\textbf{Four unknowns}: $v$, $\eta$, $C_0$, $\mu$
	
	\textbf{Four equations}:
	\begin{align*}
		v &= \left(1 - \frac{1}{\lambda}\right) C_0 - R(\eta) + \beta(1 + \eta - \mu) v \\
		R'(\eta) &= \beta v \\
		v &= c_e \\
		\frac{C_0}{\lambda} + R(\eta) + \nu c_e &= L \qquad (\nu = \mu-\eta)
	\end{align*}
	
	\bigskip
	\textbf{Growth rate}: $\gamma = \mu \log(\lambda)$
	
	\begin{itemize}
		\item Growth driven by creative destruction rate $\mu$ and innovation step $\lambda$
		\item Both incumbent innovation $\eta$ and entry $\nu$ contribute to growth
	\end{itemize}
\end{frame}

\begin{frame}{Firm Dynamics and Gibrat's Law}
	\textbf{Firm size} = number of product lines
	
	\bigskip
	\textbf{Average growth rate of the firm size (number of product lines per firm)}: $\eta - \mu = -\nu < 0$
	\begin{itemize}
		\item Net growth rate is \textbf{negative} (because entry displaces incumbents)
		\item Same for all firms regardless of size $\Rightarrow$ \textbf{Gibrat's Law}
	\end{itemize}
	
	\bigskip
	\textbf{Advantage over exogenous-shock models}:
	\begin{itemize}
		\item Policies affect the productivity process itself
		\item Can analyze how innovation incentives shape firm dynamics
		\item Clear firm definition through product lines
	\end{itemize}
\end{frame}

\begin{frame}{Generating a Pareto Firm-Size Distribution}
	\textbf{Problem}: Basic Klette--Kortum doesn't produce a Pareto tail (average growth is negative).
	
	\bigskip
	\textbf{Fix}: Suppose large firms have positive constant growth rate $g > 0$, with exit probability $\delta$. When the density at size $n$ is $h(n)$,
	
	\medskip
	Stationarity requires: $(1+g) h((1+g)n) = (1-\delta) h(n)$
	
	\medskip
	Guess Pareto distribution: $h(n) = F n^{-(\zeta + 1)}$ $\Rightarrow$
	\[
 {\zeta = -\frac{\log(1-\delta)}{\log(1+g)}}
	\]
	
	\begin{itemize}
		\item $\zeta$ small (thick tail) when $\delta$ small or $g$ large
		\item \textbf{Modification of the Klette-Kortum model to have $g>0$}: Allow some innovation to create \textbf{new} products.
		\item If new product creation $\xi$ is large enough: $g=\eta - \mu = \xi - \nu > 0$
	\end{itemize}
\end{frame}

\begin{frame}{Counterfactual Predictions and Extensions}
	\textbf{Limitations of the basic model}:
	\begin{enumerate}
		\item All product lines earn equal profits
		\begin{itemize}
			\item Unit elasticity $\Rightarrow$ revenue independent of quality
			\item Fix: higher substitutability ($\sigma > 1$)
		\end{itemize}
		
		\bigskip
		\item No cumulative innovation advantage for incumbents
		\begin{itemize}
			\item Any firm can innovate at the same cost
			\item Fix: allow incumbents to improve own products
		\end{itemize}
		
		\bigskip
		\item No Pareto tail in firm-size distribution
		\begin{itemize}
			\item Fix: allow new product creation alongside creative destruction
			\item Luttmer (2011) discusses related insights
		\end{itemize}
	\end{enumerate}
\end{frame}

% ============================================================
\section{Summary}
% ============================================================

\begin{frame}{Key Takeaways}
	\begin{enumerate}
		\item \textbf{Aggregation matters}: Firm heterogeneity affects aggregate productivity through endogenous resource allocation
		
		\item \textbf{Data}: U.S.\ shows dispersed firm sizes, large reallocation, declining dynamism, rising concentration and markups
		
		\item \textbf{Misallocation}: Firm-specific distortions reduce aggregate TFP; dispersion in distortions (not level) is what matters
		
		\item \textbf{General equilibrium}: Firing taxes and entry barriers have substantial effects on output, productivity, and reallocation
		
		\item \textbf{Market power}: Monopolistic competition gives constant markups; oligopoly (Atkeson-Burstein) generates endogenous markups linked to market shares
		
		\item \textbf{Business cycles}: Fat tails and production networks can amplify firm-level shocks to aggregate fluctuations
		
		\item \textbf{Endogenous productivity}: Klette-Kortum links innovation, creative destruction, and firm dynamics
	\end{enumerate}
\end{frame}




\end{document}

